\uuid{IrfH}
\exo7id{4150}
\titre{exo7 4150}
\auteur{quercia}
\organisation{exo7}
\datecreate{2010-03-11}
\isIndication{false}
\isCorrection{false}
\chapitre{Fonction de plusieurs variables}
\sousChapitre{Dérivée partielle}
\cpge{mpsi-mp2i.fonctions-deux-variables.c, pcsi.fonctions-deux-variables.c, ptsi.fonctions-deux-variables.c, mp-mpi.calcul-differentiel-optimisation.b, pc.fonctions-plusieurs-variables.c, psi.fonctions-plusieurs-variables.c}
\module{Analyse}
\niveau{L2}
\difficulte{}

\contenu{
\texte{
Soit $\varphi : {\R^2} \to {\R^2}$ une application affine.
}
\begin{enumerate}
    \item \question{Montrer que la matrice jacobienne, $J$, de $\varphi$ est constante.}
    \item \question{Soit $f : {\R^2} \to \R$ de classe $\mathcal{C}^2$. Pour $A \in \R^2$, on note
    $H_f(A) = \begin{pmatrix} \frac{\partial^2 f}{\partial x^2}(A) &\frac{\partial^2 f}{\partial y \partial x}(A) \cr \frac{\partial^2 f}{\partial x \partial y}(A) &\frac{\partial^2 f}{\partial y^2}(A) \cr\end{pmatrix}$
    ({\it matrice Hessienne de $f$}). Montrer que :
    $\forall\ A \in \R^2,\ H_{f\circ\varphi}(A) = {}^tJH_f(\varphi(A))J$.}
\end{enumerate}
}
